\subsection{Виртуальные перемещения}
%%%%%%%%%%%%%%%%%%%%%%%%%%%%%%%%%%%%%%%%%%%%%%%%%%%%
%%%%%%%%%%%%%%%%%%%%%%%%%%%%%%%%%%%%%%%%%%%%%%%%%%%%
%%%%%%%%%%%%%%%%%%%%%%%%%%%%%%%%%%%%%%%%%%%%%%%%%%%%
\z[{\begin{wrapfigure}[3]{r}{0.2\textwidth}%
		\vspace{-35pt}%
		\begin{flushright}%
 			\includegraphics[width=0.2\textwidth]{6.4.1}%
		\end{flushright}%
	\end{wrapfigure}%
	}]%
{Цепочка массой $m$ и длиной $l$ подвешена так, что расстояние между концами (по высоте) равно $h$. Найдите разность сил натяжения на концах цепочки.}
%
{$\Delta T =mg\frac{h}{l}$}

%%%%%%%%%%%%%%%%%%%%%%%%%%%%%%%%%%%%%%%%%%%%%%%%%%%%
\z{Нить, перекинутая через блок, изначально не натянута. К свободному концу нити прикладывают силу $F$. На сколько сместится конец нити под действием этой силы? Определите эквивалентный коэффициент жесткости.
		\begin{center}%
 			\includegraphics[width=0.62\textwidth]{6.4.2}%
		\end{center}%
}
%
{\puan $\frac{k}{5}$;~\puan $\frac{4k}{17}$;~\puan $\frac{k}{5}$;~\puan $\frac{k}{6}$}

%%%%%%%%%%%%%%%%%%%%%%%%%%%%%%%%%%%%%%%%%%%%%%%%%%%%
\z{Найдите эффективную жёсткость системы. Внешняя сила прикладывается к ней в месте, обозначенном на рисунке стрелкой.
		\begin{center}%
 			\includegraphics[width=0.62\textwidth]{6.4.3}%
		\end{center}%
		}
%
{\puan $k$;~\puan $\frac{2k}{51}$}


%%%%%%%%%%%%%%%%%%%%%%%%%%%%%%%%%%%%%%%%%%%%%%%%%%%%
\z[{\begin{wrapfigure}[3]{r}{0.3\textwidth}%
		\vspace{-10pt}%
		\begin{flushright}%
 			\includegraphics[width=0.3\textwidth]{6.4.4}%
		\end{flushright}%
	\end{wrapfigure}%
	}]%
{На стержне в точках $A$ и $B$ подвешены грузы массой $0.2 \, \text{кг}$ и $0.4 \, \text{кг}$ соответственно. Определите, какую силу надо приложить в точке $C$, чтобы удержать систему в равновесии. $OA = 0.4 \, \text{м}$, $OB = 0.1 \, \text{м}$, $BC = 0.3 \, \text{м}$.}
%
{$F = 1 \, \text{Н}$}


%%%%%%%%%%%%%%%%%%%%%%%%%%%%%%%%%%%%%%%%%%%%%%%%%%%%
\z[{\begin{wrapfigure}[3]{r}{0.15\textwidth}%
		\vspace{-20pt}%
		\begin{flushright}%
 			\includegraphics[width=0.15\textwidth]{6.4.5}%
		\end{flushright}%
	\end{wrapfigure}%
	}]%
{В системе, изображённой на рисунке, к нижнему шарниру подвешен груз массы $m$. Определите, какую силу $F$ нужно приложить к свободному концу нити, чтобы система оставалась в равновесии.}
%
{$F = 2mg$}

%%%%%%%%%%%%%%%%%%%%%%%%%%%%%%%%%%%%%%%%%%%%%%%%%%%%
\z[{\begin{wrapfigure}[5]{r}{0.3\textwidth}%
		\vspace{-10pt}%
		\begin{flushright}%
 			\includegraphics[width=0.3\textwidth]{6.4.6}%
		\end{flushright}%
	\end{wrapfigure}%
	}]%
{Из лёгких жёстких стержней, соединённых шарнирно и образующих три одинаковых ромба, и лёгкого рычага собрана система, к которой подвешены два груза. Определите силу натяжения нити $T$, удерживающей систему в равновесии, если масса $m = 1 \ \text{кг}$.}
%
{$T = \frac{16}{3} mg$}

%%%%%%%%%%%%%%%%%%%%%%%%%%%%%%%%%%%%%%%%%%%%%%%%%%%%
\z[{\begin{wrapfigure}[4]{r}{0.3\textwidth}%
		\vspace{-0pt}%
		\begin{flushright}%
 			\includegraphics[width=0.15\textwidth]{6.4.7}%
		\end{flushright}%
	\end{wrapfigure}%
	}]%
{На гладком закреплённом шкиве радиусом $R$ висит массивный однородный канат массой $m$ и длиной $l = 8R$ (см. рисунок).  Ось $O$ шкива горизонтальна. Найти силу натяжения каната в точке $A$.}
%
{$T = \frac{8 - \pi}{16} mg$}

